\documentclass[11pt]{amsart}
\usepackage{amssymb}
\usepackage[margin=1.15in]{geometry}
\usepackage{array}
\usepackage{longtable}
\usepackage[hidelinks]{hyperref}
\hypersetup{pdftitle={Jacobian localization and regular sequences in two variables}}

\newtheorem{theorem}{Theorem}[section]
\newtheorem{proposition}[theorem]{Proposition}
\newtheorem{lemma}[theorem]{Lemma}
\newtheorem{corollary}[theorem]{Corollary}
\theoremstyle{definition}
\newtheorem{definition}[theorem]{Definition}
\theoremstyle{remark}
\newtheorem{remark}[theorem]{Remark}
\DeclareMathOperator{\Ann}{Ann}
\newcommand{\content}{\mathrm{c}}
\newcommand{\code}[1]{\texttt{#1}}
\newcommand{\mere}[1]{\left\lVert #1\right\rVert}

\title[Jacobian localization and regular sequences]{Jacobian localization and regular sequences\\in two variables}
\date{October 25, 2026}

% stamp on every page: file name and date/time of compilation
\newcount\stamphour \newcount\stampmin
\stamphour=\time \divide\stamphour by 60
\stampmin=\time \multiply\stamphour by 60 \advance\stampmin by -\stamphour \divide\stamphour by 60
\edef\stamptime{\the\year-\ifnum\month<10 0\fi\the\month-\ifnum\day<10 0\fi\the\day\space
 \ifnum\stamphour<10 0\fi\the\stamphour:\ifnum\stampmin<10 0\fi\the\stampmin}
\AddToHook{shipout/foreground}{\put(0.5\paperwidth,-\dimexpr\paperheight-0.45in\relax)%
 {\makebox[0pt]{\scriptsize\ttfamily\jobname.pdf\quad\stamptime}}}
\begin{document}
\begin{abstract}
For polynomials $F,G$ over any commutative ring, their Jacobian
determinant $J$ has the following property: $FU=GV$ implies
$J^eU=GW$ for some $e\geq0$ and some polynomial $W$.
We prove this by a finite Taylor lifting construction and stabilization
of annihilators, after reduction to a noetherian coefficient ring.
Together with a content argument, this proves that $(F,G)$ is regular
in both orders after inverting $J$. The proof uses no flatness theorem.
The final section describes the Unit formalisation, which proves the
divisibility assertion under an explicit annihilator-stabilization
hypothesis; the noetherian reduction remains outside the formalisation.
\end{abstract}
\maketitle

\section{Statement}

All rings are commutative with identity; the zero ring is allowed.
Let $A$ be a ring, let $R=A[X,Y]$, and let $F,G\in R$. Write
\[
 J=F_XG_Y-F_YG_X,\qquad S=R[1/J],
\]
where subscripts denote formal partial derivatives.
An element is \emph{regular} if multiplication by it is injective.
We call $(a,b)$ \emph{commutatively regular} if $a,b$ are regular
individually, $b$ is regular modulo $(a)$, and $a$ is regular modulo
$(b)$. Here regularity of a sequence does not require its generated
ideal to be proper.

\begin{theorem}\label{thm:main}
The pair $(F,G)$ is commutatively regular in $S$.
Moreover, for all $U,V\in R$,
\begin{equation}\label{eq:divisibility}
 FU=GV\quad\Longrightarrow\quad
 \exists e\geq0\;\exists W\in R\quad J^eU=GW.
\end{equation}
\end{theorem}

The individual regularity of $F$ and $G$ is Step~I. The principal
point is Step~II, namely~\eqref{eq:divisibility}. We first give a
conditional proof of Step~II using only finite polynomial calculations,
then supply its hypothesis by a noetherian reduction.

\section{Step I: content and individual regularity}\label{sec:content}

For a polynomial $H$ in finitely many indeterminates over $A$, let
$\content(H)\subseteq A$ denote the ideal generated by its coefficients.
The following argument is the one in the supplied note
\code{CAR1.pdf}~\cite{CAR1}.

\begin{lemma}\label{lem:content}
Let $H,P\in A[\mathbf X]$. If
$\content(P)\subseteq\content(H)$, then $H$ is regular in
$A[\mathbf X][1/P]$.
\end{lemma}
\begin{proof}
We use the consequence of Dedekind--Mertens that, for polynomials
$H,Q$, there is an integer $m$ such that
\[
 \content(H)^m\content(Q)\subseteq\content(HQ).
\]
Suppose a polynomial $K$ satisfies $HK=0$ after localization.
Then $P^kHK=0$ in $A[\mathbf X]$ for some $k\geq0$.
Apply the displayed inclusion to $Q=P^kK$. For a suitable $m$,
\[
 \content(P^{m+k}K)
 \subseteq \content(P)^m\content(P^kK)
 \subseteq \content(H)^m\content(P^kK)
 \subseteq \content(HP^kK)=0.
\]
Hence $P^{m+k}K=0$, so $K=0$ after localization.
Every localized element has a polynomial numerator, which proves
the assertion.
\end{proof}

Every coefficient of a partial derivative of $F$ belongs to
$\content(F)$, and similarly for $G$. Consequently
\[
 \content(J)\subseteq\content(F),\qquad
 \content(J)\subseteq\content(G).
\]
Lemma~\ref{lem:content}, with $P=J$, proves the individual regularity
of $F$ and $G$ in $S$. Notice the direction of the content inclusion
in the lemma.

\section{Two finite algebraic lemmas}

\begin{definition}\label{def:stable}
For a ring $B$, an element $f\in B$ and an integer $N\geq0$, write
$\operatorname{Stable}_B(f,N)$ for the condition
\begin{equation}\label{eq:stable}
 \forall v\in B,\qquad f^{N+1}v=0\ \Longrightarrow\ f^Nv=0.
\end{equation}
Equivalently, $\Ann_B(f^N)=\Ann_B(f^{N+1})$.
\end{definition}

\begin{lemma}[A truncated linear equation]\label{lem:linear}
Assume $\operatorname{Stable}_B(f,N)$. If
\[
 (f-T)p(T)=0\pmod {T^{N+1}},
\]
then $p(0)=0$.
\end{lemma}
\begin{proof}
Write $p(T)=u_0+u_1T+\cdots+u_NT^N$ modulo $T^{N+1}$.
Comparison of coefficients gives
\[
 fu_0=0,\qquad fu_i=u_{i-1}\quad(1\leq i\leq N).
\]
Thus $u_0=f^Nu_N$ and $f^{N+1}u_N=0$.
Condition~\eqref{eq:stable} implies $f^Nu_N=0$, hence $u_0=0$.
This also covers $N=0$: the hypothesis then says that $f$ itself
is regular.
\end{proof}

\begin{lemma}[Finite Taylor lifting]\label{lem:taylor}
Let $B$ be an $A$-algebra, and let $x,y\in B$ satisfy
\[
 G(x,y)=0,\qquad J(x,y)\in B^\times.
\]
Put $f=F(x,y)$. For every $N\geq0$, there are polynomials
$x_N(T),y_N(T)\in B[T]$ such that
\begin{align}
 x_N(0)&=x,& y_N(0)&=y,\label{eq:initial}\\
 F(x_N(T),y_N(T))&\equiv f-T\pmod {T^{N+1}},&
 G(x_N(T),y_N(T))&\equiv0\pmod {T^{N+1}}.
 \label{eq:lift}
\end{align}
\end{lemma}
\begin{proof}
Start with the constant polynomials $x,y$, which give the assertions
modulo $T$. Suppose $x_*,y_*$ have already been constructed modulo
$T^r$, where $r\geq1$, and write their next errors as
\[
 F(x_*,y_*)-(f-T)\equiv\epsilon_FT^r,\qquad
 G(x_*,y_*)\equiv\epsilon_GT^r
 \pmod {T^{r+1}}.
\]
We seek corrections $x_*+aT^r$ and $y_*+bT^r$.
For every polynomial $H$, the ordinary polynomial expansion gives
\begin{equation}\label{eq:taylor-step}
 H(x_*+aT^r,y_*+bT^r)
 \equiv H(x_*,y_*)+
 T^r\bigl(H_X(x,y)a+H_Y(x,y)b\bigr)
 \pmod {T^{r+1}}.
\end{equation}
Indeed, terms containing at least two corrections are divisible by
$T^{2r}$, hence by $T^{r+1}$; in the linear terms, the derivatives
may be evaluated at $x,y$ since $x_*,y_*$ have those constant terms.
This argument is valid in every characteristic.

The matrix
\[
 M=\begin{pmatrix}
 F_X(x,y)&F_Y(x,y)\\
 G_X(x,y)&G_Y(x,y)
 \end{pmatrix}
\]
has determinant $J(x,y)$, a unit. Therefore we can take
\[
 \binom ab=-M^{-1}\binom{\epsilon_F}{\epsilon_G},\qquad
 M^{-1}=J(x,y)^{-1}
 \begin{pmatrix}G_Y(x,y)&-F_Y(x,y)\\-G_X(x,y)&F_X(x,y)\end{pmatrix}.
\]
By~\eqref{eq:taylor-step}, both errors then vanish modulo $T^{r+1}$.
The constant terms are unchanged. Performing these corrections for
$r=1,\ldots,N$ proves the lemma. For $N=0$ the initial constants suffice.
\end{proof}

The construction uses only finitely many polynomial operations and the
inverse of $J(x,y)$. In particular, it does not require division by
factorials or the construction of an infinite formal power series.

\section{Step II under the stabilization hypothesis}\label{sec:conditional}

Let
\[
 B=\bigl(R/(G)\bigr)[1/J],
\]
where the localization is at the image of $J$. Denote the images of
$X,Y,F,U$ in $B$ by $x,y,f,u$, respectively.

\begin{proposition}\label{prop:conditional}
Suppose there exists $N\geq0$ with
$\operatorname{Stable}_B(f,N)$. Then~\eqref{eq:divisibility} holds
for every $U,V\in R$.
\end{proposition}
\begin{proof}
Assume $FU=GV$. In $B$ we have $G(x,y)=0$ and $J(x,y)$ is a unit,
so Lemma~\ref{lem:taylor} supplies $x_N,y_N$.
Substitution into the polynomial identity $FU=GV$ gives
\[
 (f-T)U(x_N(T),y_N(T))\equiv0\pmod {T^{N+1}}.
\]
Lemma~\ref{lem:linear} implies
\[
 u=U(x_N(0),y_N(0))=0\quad\text{in }B.
\]
The kernel criterion for localization now gives $e\geq0$ such that
$J^eU=0$ in $R/(G)$. Thus $J^eU\in(G)$ in $R$, which is precisely
$J^eU=GW$ for some $W\in R$.
\end{proof}

\section{Noetherian reduction and completion of the proof}\label{sec:noetherian}

\begin{proof}[Proof of \eqref{eq:divisibility} over an arbitrary ring]
Fix a relation $FU=GV$. Let $A_0\subseteq A$ be the subring generated
by the coefficients of these four polynomials and by $1$.
It is a finitely generated $\mathbb Z$-algebra, hence noetherian by
the Hilbert basis theorem. Since $A_0$ is a subring, the given
identity already holds in $A_0[X,Y]$.

Set
\[
 B_0=\bigl(A_0[X,Y]/(G)\bigr)[1/J],
 \qquad f_0=\text{the image of }F\text{ in }B_0.
\]
Polynomial extensions, quotients and localizations of noetherian rings
are noetherian. Hence the ascending chain
\[
 \Ann_{B_0}(1)\subseteq\Ann_{B_0}(f_0)
 \subseteq\Ann_{B_0}(f_0^2)\subseteq\cdots
\]
stabilizes. In particular, $\operatorname{Stable}_{B_0}(f_0,N)$ holds
for some $N$. Proposition~\ref{prop:conditional}, applied over $A_0$,
gives $e\geq0$ and $W\in A_0[X,Y]$ with $J^eU=GW$.
This identity also holds in $A[X,Y]$.
\end{proof}

This reduction supplies the stabilization hypothesis in $B_0$ for
the particular relation under consideration. It does not assert that
the ring $B$ formed over the original arbitrary $A$ is noetherian.
The finite argument of Section~\ref{sec:conditional} is constructive
under its stated hypothesis; the reduction just given uses classical
noetherian algebra.

\begin{proof}[Completion of Theorem~\ref{thm:main}]
Step~I gives the individual regularity of $F$ and $G$ in $S$.
To see that $F$ is regular modulo $G$, work in
$S/(G)\cong(R/(G))[1/J]$. Suppose an element represented by $U/J^r$
is killed by $F$. Clearing the localization equality gives
\[
 J^kFU=GV
\]
in $R$ for some $k\geq0$ and $V\in R$.
Apply~\eqref{eq:divisibility} to $F(J^kU)=GV$. We obtain
$J^{e+k}U\in(G)$ for some $e$, so $U/J^r=0$ in $S/(G)$.

Exchanging $F$ and $G$ changes the Jacobian to $-J$, which defines
the same localization. The same argument therefore proves that $G$
is regular modulo $F$. This proves regularity in both orders.
\end{proof}

\section{What has been formalised in Unit}\label{sec:formalisation}

The development is in \code{unit/sag/jacobian/}. The theorem checked
by Unit is Proposition~\ref{prop:conditional}, including the
construction of $B$ and the final certificate in $A[X,Y]$.
Its only additional algebraic hypothesis is
\begin{equation}\label{eq:formal-input}
 \mere{\sum_{N:\mathbb N}\operatorname{Stable}_B(f,N)}.
\end{equation}
Here $\mere{\cdot}$ denotes propositional truncation, or mere existence.
The public theorem takes an actual polynomial identity $FU=GV$ and
returns
\begin{equation}\label{eq:formal-output}
 \mere{\sum_{e:\mathbb N}\sum_{W:A[X,Y]}(J^eU=GW)}.
\end{equation}
Thus neither $N$ nor the output witnesses need be chosen globally.
The untruncated type inside~\eqref{eq:formal-input} is named
\code{NoetherianInput}; its name does not assert a formalised
noetherianity theorem.

\subsection*{Polynomial representation and the public interface}

The source ring is the existing coefficient-quotient polynomial ring
$A[X][Y]$. The inputs $F,G$ are finite polynomial expressions with
constants in $A$, two variables, sums and products. Partial derivatives
are evaluated by the usual formal sum and product rules.
\code{ExpressionRepresentation} proves that every element of
$A[X][Y]$ merely has such an expression. The inputs $U,V$ to the
public theorem are arbitrary elements of this polynomial ring.
All identity hypotheses are equalities in that ring, rather than
equalities of polynomial functions.

The public interface, with the foundation parameters understood, is
\begin{verbatim}
module D = JacobianDivisibility[l].Over[A,setA,F,G]
D.result input u v relation
\end{verbatim}
Here \code{setA} expresses that the carrier of the commutative ring is
a set, \code{input} has type~\eqref{eq:formal-input}, and
\code{relation} is the equality $FU=GV$.

\subsection*{Correspondence with the proof}

The following table lists the source files in the development directory.
Each file name has extension \code{.seg}.

\begin{longtable}{@{}p{0.39\textwidth}p{\dimexpr0.61\textwidth-2\tabcolsep\relax}@{}}
\textbf{File or files}&\textbf{Checked argument}\\[4pt]
\endhead
\code{FiniteJets}&
Equality of coefficients below a finite degree, and the product rule
for changing a coefficient of positive degree.\\[5pt]
\code{StableAnnihilator}&
Lemma~\ref{lem:linear}, including $N=0$.\\[5pt]
\code{Bivariate},\newline\code{BivariateNaturality}&
Polynomial expressions, evaluation and evaluated partial derivatives;
Taylor's correction formula and compatibility with ring maps.\\[5pt]
\code{LinearSolve},\newline\code{TaylorLift}&
Inversion of the $2\times2$ Jacobian matrix and the finite construction
of Lemma~\ref{lem:taylor}.\\[5pt]
\code{PolynomialIdentity},\newline\code{RegularAt}&
Substitution of the genuine polynomial identity into the finite lift,
and the conclusion $u=0$ in the target ring.\\[5pt]
\code{PrincipalQuotient},\newline\code{PowerLocalization}&
The quotient by $G$ and localization at $J$, with their proved kernel
certificates. The localization reuses the existing fraction library.\\[5pt]
\code{ExpressionRepresentation}&
Expression representatives for all elements of $A[X][Y]$, by
quotient and finite list induction.\\[5pt]
\code{JacobianDivisibility}&
Proposition~\ref{prop:conditional} and the conclusion
\eqref{eq:formal-output}, after clearing denominators.\\[5pt]
\code{Examples}&
For $F=X+X^2$, $G=Y$ at $(0,0)$, all finite lifts and the coefficients
$x(T)=-T-T^2\bmod T^3$, $y(T)=0\bmod T^3$, over any commutative ring.
\end{longtable}

\subsection*{Foundations, validation and remaining work}

The construction uses the same univalence and replacement parameters
as the existing Unit polynomial library, from which the effective set
quotients and propositional truncations are obtained. It introduces
no excluded-middle or choice axiom. All recursive definitions use
Unit's data eliminators.
\code{JetIdentities.seg} contains ordinary equality proofs generated
by the existing ring-identity generator. Unit checks these proof terms;
the generator is not an additional trusted principle.

From \code{unit/}, the complete checking command is
\begin{verbatim}
./segc sag/jacobian all
\end{verbatim}
The check on October 25, 2026 reported \emph{40 modules, 0 failed}.
This count includes the imported library modules.

The remaining unformalised part is the derivation of
\eqref{eq:formal-input} by the noetherian argument in
Section~\ref{sec:noetherian}: the finite coefficient reduction, the
Hilbert basis theorem and preservation of noetherianity, and stabilization
of the annihilator chain. Step~I from Section~\ref{sec:content} is also
not formalised in this directory. The separate development
\code{unit/sag/mertens/} now proves the univariate Dedekind--Mertens
formula over every commutative ring, following the supplied direct
proof~\cite{Mertens}. It constructs signed-sum certificates for the
stronger additive-subgroup inclusion, proves the reverse inclusion,
and obtains both ideal inclusions. Its check reports
\emph{14 modules, 0 failed}. The multivariate content lemma and
localization application used in Step~I have not yet been connected
to that result. A separate note, \code{dedekind-mertens.tex}, presents
the proof and its formalisation status.

Consequently the unrestricted
Theorem~\ref{thm:main} is proved in this note, while the checked Unit
result at present is the conditional Step~II described above.

\begin{thebibliography}{1}
\bibitem{CAR1}
\emph{ContentAndRegularity.tex}, one-page note supplied as
\code{CAR1.pdf}, compilation dated October 25, 2026.
\bibitem{Mertens}
\emph{A direct proof of the Dedekind-Mertens Lemma}, supplied as
\code{mertens.pdf}, April 5, 2006.
\end{thebibliography}

\end{document}
